\documentclass[14pt]{extarticle}
\usepackage{fontspec}
\usepackage{polyglossia}
\usepackage{amsmath}
\usepackage{amssymb}
\usepackage{xcolor}
\usepackage{fancyhdr}
\usepackage{graphicx}
\usepackage{listings}
\usepackage[most]{tcolorbox}
% \usepackage{setspace}
\usepackage{pifont}
\usepackage{enumitem}
\usepackage{titlesec}
\usepackage[bottom]{footmisc}
\usepackage{titling}
\usepackage{minted}
\usepackage{etoolbox}
\usepackage{array}
\usepackage{longtable} % page-breakable tables (must load before polyglossia/bidi RTL setup)
\usepackage{extsizes}
\usepackage{float}
\usepackage{booktabs}
\usepackage{hyperref}
\usepackage[onehalfspacing]{setspace}
\newif\ifwhiteblack
\ifdefined\setwhiteblack
  \setwhiteblack
\fi

\ifwhiteblack
  \usepackage[a4paper, landscape, top=1.25cm, bottom=1.25cm, left=1.25cm, right=1.25cm, headheight=1.5cm, headsep=0.5cm]{geometry}
\else
  \usepackage[a4paper, top=2.7cm, bottom=2.7cm, left=2.5cm, right=2.5cm, headheight=1.5cm, headsep=0.5cm]{geometry}
\fi
\usepackage{tikz}
\usepackage{datetime2}

\usetikzlibrary{automata, positioning, arrows.meta, arrows}
\providecommand{\LTcaptype}{table}

\newfontfamily\emoji{Segoe UI Emoji}

\pagestyle{fancy}

\setmainlanguage[numerals=western]{arabic}
\setotherlanguage{english}
\newfontfamily\arabicfont[Script=Arabic]{Calibri}
\newfontfamily\arabicfonttt[Script=Arabic]{Courier New}



\ifwhiteblack
\lstset{
  language=[Sharp]C,
  numbers=left,
  stepnumber=1,
  numbersep=8pt,
  frame=single,
  basicstyle=\ttfamily\small,
  keywordstyle=\bfseries,
  stringstyle=\ttfamily,
  commentstyle=\itshape
}
\else
\lstset{
  language=[Sharp]C,
  numbers=left,
  stepnumber=1,
  numbersep=8pt,
  frame=single,
  basicstyle=\ttfamily\small,
  keywordstyle=\color{blue},
  stringstyle=\color{red},
  commentstyle=\color{green!50!black}
}
\fi

\newif\ifdetailed
\ifdefined\setdetailed
  \setdetailed
\fi

\newif\ifwithsols
\ifdefined\setwithsols
  \setwithsols
\fi

\ifwhiteblack
  \AtBeginDocument{
    \LARGE
    \usemintedstyle{bw}
  }
\fi

% Default minted look: compact code that wraps long lines.
% Not applied in black-and-white landscape mode, which already uses \LARGE.
\ifwhiteblack\else
  \setminted{fontsize=\small, breaklines=true}
\fi

% unified tcolorboxes for articles
\ifwhiteblack
\tcbset{colback=white, colframe=black, fonttitle=\bfseries, boxrule=1pt}
\newtcolorbox{boxDef}[1][]{title={{\emoji📘} تعريف\ifx\\#1\\\else ~#1\fi :}}
\newtcolorbox{boxExercise}[1][]{title={{\emoji🧩} تمرين\ifx\\#1\\\else ~#1\fi :}}
\newtcolorbox{boxExample}[1][]{title={{\emoji📝} مثال\ifx\\#1\\\else ~#1\fi :}}
\newtcolorbox{boxNote}[1][]{title={{\emoji✨} ملاحظة\ifx\\#1\\\else ~#1\fi :}}
\newtcolorbox{boxAttention}[1][]{title={{\emoji🔔} تنبيه\ifx\\#1\\\else ~#1\fi :}}
\newtcolorbox{boxWarning}[1][]{title={{\emoji⚡} ملاحظة هامة\ifx\\#1\\\else ~#1\fi :}}
\newtcolorbox{boxSolution}[1][]{title={{\emoji✅} حل\ifx\\#1\\\else ~#1\fi :}}
\newtcolorbox{boxSymbol}[1][]{title={{\emoji🔣} رمز\ifx\\#1\\\else ~#1\fi :}}
\newtcolorbox{boxHint}[1][]{title={{\emoji💡} تلميح\ifx\\#1\\\else ~#1\fi :}}
\else
\tcbset{colback=white, colframe=black, fonttitle=\bfseries, boxrule=0.8pt}
\newtcolorbox{boxDef}[1][]{colback=blue!5!white,colframe=blue!75!black,
  title={{\emoji📘} تعريف\ifx\\#1\\\else ~#1\fi :}}
\newtcolorbox{boxExercise}[1][]{colback=cyan!5!white,colframe=cyan!70!black,
  title={{\emoji🧩} تمرين\ifx\\#1\\\else ~#1\fi :}}
\newtcolorbox{boxExample}[1][]{colback=yellow!5!white,colframe=orange!90!black,
  title={{\emoji📝} مثال\ifx\\#1\\\else ~#1\fi :}}
\newtcolorbox{boxNote}[1][]{colback=gray!10!white,colframe=black,
  title={{\emoji✨} ملاحظة\ifx\\#1\\\else ~#1\fi :}}
\newtcolorbox{boxAttention}[1][]{colback=magenta!10!white,colframe=magenta!80!black,
  title={{\emoji🔔} تنبيه\ifx\\#1\\\else ~#1\fi :}}
\newtcolorbox{boxWarning}[1][]{colback=red!5!white,colframe=red!75!black,
  title={{\emoji⚡} ملاحظة هامة\ifx\\#1\\\else ~#1\fi :}}
\newtcolorbox{boxSolution}[1][]{colback=green!5!white,colframe=green!60!black,
  title={{\emoji✅} حل\ifx\\#1\\\else ~#1\fi :}}
\newtcolorbox{boxSymbol}[1][]{colback=purple!5!white,colframe=purple!70!black,
  title={{\emoji🔣} رمز\ifx\\#1\\\else ~#1\fi :}}
\newtcolorbox{boxHint}[1][]{colback=teal!5!white,colframe=teal!60!black,
  title={{\emoji💡} تلميح\ifx\\#1\\\else ~#1\fi :}}
\fi


\ifwhiteblack
\tcbset{simplecode/.style={ colback=white, colframe=black, boxrule=1pt, arc=2pt, left=4pt,right=4pt,top=4pt,bottom=4pt}}
\else
\tcbset{simplecode/.style={ colback=gray!5, colframe=black!50, boxrule=0.4pt, arc=2pt, left=4pt,right=4pt,top=4pt,bottom=4pt}}
\fi
% Compact code boxes: tighter line spacing (the preamble uses onehalfspacing) and small vertical gaps.
% Must come AFTER the simplecode definition above.
\tcbset{simplecode/.append style={fontupper=\setstretch{0.95}, top=4pt, bottom=4pt, before skip=6pt, after skip=6pt}}
\newenvironment{boxCode}{\begin{tcolorbox}[simplecode]}{\end{tcolorbox}}

% Code cell used inside tables (keeps code left-to-right inside RTL tables).
% Use inside a fixed-width p{...} column: \begin{pycell}\begin{minted}...\end{minted}\end{pycell}
\newtcolorbox{pycell}{frame empty,colback=white,boxsep=0pt,left=2pt,right=2pt,top=2pt,bottom=2pt,width=\linewidth}

\newcolumntype{C}[1]{>{\centering\arraybackslash}p{#1}}

\makeatletter
\DTMnewdatestyle{dotted}{%
  \renewcommand{\DTMdisplaydate}[4]{\two@digits{##3}.\two@digits{##2}.##1}%
  \renewcommand{\DTMDisplaydate}{\DTMdisplaydate}%
}
\makeatother
\DTMsetdatestyle{dotted}

% redefine spaces after titles
\makeatletter
\renewcommand{\@maketitle}{%
  \begin{center}
    {\huge \bfseries \@title \par}%
    \vskip 0.2em % space between title and author
    {\large \@author \par}%
    % \vskip 0.2em % space between author and date
    % {\normalsize \@date \par}%
    {\small \hfill \begin{english}\DTMtoday\end{english} \par}%
  \end{center}
}
\makeatother

\fancyhf{} % clear default
\fancypagestyle{plain}{
  \fancyhf{}
  \fancyhead[L]{مدرسة التسامح الشاملة}
  % \fancyhead[L]{\includegraphics[height=1cm]{../../../images/logoTasamoh.png}}
  \fancyhead[R]{الأستاذ محمود اغبارية}
  \fancyfoot[C]{\thepage}
}

\fancyhead[L]{مدرسة التسامح الشاملة}
\fancyhead[R]{الأستاذ محمود اغبارية}
\fancyfoot[C]{\thepage}

\setcounter{tocdepth}{3} % only section subsection and subsubsection in TOC

\newcommand{\importpdfpage}[2]{%
    \thispagestyle{empty}
    \newgeometry{left=0cm,right=0cm,top=0cm,bottom=0cm}
    \includegraphics[page=#2,width=0.95\paperwidth,height=0.95\paperheight,keepaspectratio]{#1}
    \restoregeometry
}

\newcommand{\insertFullImg}[1]{%
    \noindent
    \makebox[\textwidth][c]{
        \includegraphics[width=0.9\paperwidth,keepaspectratio]{#1}%
    }%
}%


\ifwithsols
\title{حل ورقة تمرن بايثون العمليات (جزء أ)}
\else
\title{ورقة تمرن بايثون العمليات (جزء أ)}
\fi

\begin{document}

\maketitle
\thispagestyle{fancy}

%==============================================================
\section*{القسم الأول: قراءة العمليات وتتبّع المخرج}

\begin{enumerate}[label=\textbf{تمرين \arabic*:}, wide=0pt, leftmargin=*, itemsep=2.2em]

% ---------------------------------------------------------------- 1
\item اكتبوا \textbf{المخرج} الناتج عن استدعاء كل واحدة من العمليات في الجدول التالي.

{\small
\renewcommand{\arraystretch}{1.3}
\begin{longtable}{|c|p{0.45\textwidth}|p{0.3\textwidth}|}
\hline
\multicolumn{1}{|c|}{\textbf{البند}} & \multicolumn{1}{c|}{\textbf{مقطع البرنامج}} & \multicolumn{1}{c|}{\textbf{المخرج}} \\
\hline
\endhead
1 &
\begin{pycell}
\begin{english}
\begin{minted}[fontsize=\small]{python}
def ex1():
    num = 34
    a = num + 10
    b = num / 10
    print(b * 10 + a)

ex1()
\end{minted}
\end{english}
\end{pycell} & \rule{0pt}{1.3cm} \\
\hline
2 &
\begin{pycell}
\begin{english}
\begin{minted}[fontsize=\small]{python}
def ex2():
    a = 10
    b = 30
    c = 20
    print(a + b + c)

ex2()
\end{minted}
\end{english}
\end{pycell} & \rule{0pt}{1.3cm} \\
\hline
3 &
\begin{pycell}
\begin{english}
\begin{minted}[fontsize=\small]{python}
def ex3():
    a = 5
    b = 2
    c = 10
    print(a * b + c)

ex3()
\end{minted}
\end{english}
\end{pycell} & \rule{0pt}{1.3cm} \\
\hline
4 &
\begin{pycell}
\begin{english}
\begin{minted}[fontsize=\small]{python}
def ex4():
    a = 3
    b = 2
    c = 1
    print(a + 2 * b + 3 * c)

ex4()
\end{minted}
\end{english}
\end{pycell} & \rule{0pt}{1.3cm} \\
\hline
5 &
\begin{pycell}
\begin{english}
\begin{minted}[fontsize=\small]{python}
def ex5():
    a = 10
    b = 30
    c = 20
    print((a + b) // c)

ex5()
\end{minted}
\end{english}
\end{pycell} & \rule{0pt}{1.3cm} \\
\hline
6 &
\begin{pycell}
\begin{english}
\begin{minted}[fontsize=\small]{python}
def ex6():
    a = 10
    b = 30
    print((b + a) / (b - a))

ex6()
\end{minted}
\end{english}
\end{pycell} & \rule{0pt}{1.3cm} \\
\hline
\end{longtable}
}

\ifwithsols
\begin{boxSolution}
\begin{english}
\renewcommand{\arraystretch}{1.2}
\begin{tabular}{rcl}
1 & $\rightarrow$ & \texttt{78.0} \\
2 & $\rightarrow$ & \texttt{60} \\
3 & $\rightarrow$ & \texttt{20} \\
4 & $\rightarrow$ & \texttt{10} \\
5 & $\rightarrow$ & \texttt{2} \\
6 & $\rightarrow$ & \texttt{2.0} \\
\end{tabular}
\end{english}

لاحظوا أنّ ناتج العامل \textenglish{/} هو دائمًا من نمط \textenglish{float} (لذلك المخرج في البند 1 هو \textenglish{78.0} وليس \textenglish{78})، بينما ناتج العامل \textenglish{//} على عددين صحيحين هو عدد صحيح.
\end{boxSolution}

\fi
\filbreak

% ---------------------------------------------------------------- 2

\item أمامكم عملية \textbf{تستقبل} من المستخدم \textbf{سعر} صندوق البيتسا و\textbf{كمية} الصناديق المطلوبة. تحسب العملية وتطبع المبلغ \textbf{الكلي} الذي يجب دفعه.

أكملوا الناقص في مقطع البرنامج التالي باستخدام مخزن الكلمات.

\begin{boxCode}
\begin{english}
\begin{minted}{python}
def calculate_pizza_order():
    # Read input from the user
    num_pizzas = ______________________________
    ________________ = float(input("Enter the price per pizza: "))

    # Calculate the total price
    total_price = num_pizzas * price_per_pizza

    # Display the result
    print("The total price for pizzas is:", ____________, "nis")

# Call the function
calculate_pizza_order()
\end{minted}
\end{english}
\end{boxCode}

\textbf{مخزن الكلمات:}
\begin{center}
\begin{english}
\begin{boxCode}
\texttt{int(input("Enter the number of pizzas: "))} \quad , \quad \texttt{price\_per\_pizza} \quad , \quad \texttt{total\_price}
\end{boxCode}
\end{english}
\end{center}

\ifwithsols
\begin{boxSolution}
\begin{english}
\begin{minted}{python}
def calculate_pizza_order():
    # Read input from the user
    num_pizzas = int(input("Enter the number of pizzas: "))
    price_per_pizza = float(input("Enter the price per pizza: "))

    # Calculate the total price
    total_price = num_pizzas * price_per_pizza

    # Display the result
    print("The total price for pizzas is:", total_price, "nis")

# Call the function
calculate_pizza_order()
\end{minted}
\end{english}
\end{boxSolution}

\fi

\filbreak
% ---------------------------------------------------------------- 3
\item أمامكم عملية \textbf{تطبع} معلومات عن المدرسة:

\begin{boxCode}
\begin{english}
\begin{minted}{python}
def school_info():
    school_name = "ABC High School"
    address = "123 Main St, Anytown"
    principal = "Mrs. Smith"
    print("School:", school_name)
    print("Address:", address)
    print("Principal:", principal)

# Call the function
school_info()
\end{minted}
\end{english}
\end{boxCode}

\begin{enumerate}[leftmargin=*, itemsep=1.2em]
\item[\textbf{أ.}] اختاروا الإجابة \textbf{الأكثر صحة}: لكي \textbf{نستقبل} اسم المدرسة من المستخدم علينا كتابة الأمر التالي:
\begin{boxCode}
\begin{english}
\begin{minted}{text}
a. input("Enter school name: ")
b. int(input("Enter school name: "))
c. float(input("Enter school name: "))
d. print("Enter school name: ")
\end{minted}
\end{english}
\end{boxCode}

\item[\textbf{ب.}] اختاروا الإجابة \textbf{الأكثر صحة}: اسم العملية هو:
\begin{boxCode}
\begin{english}
\begin{minted}{text}
a. school_info
b. principal
c. school_name
d. def
\end{minted}
\end{english}
\end{boxCode}

\item[\textbf{ج.}] عدّلوا العملية بحيث \textbf{تستقبل} من المستخدم المعلومات (اسم المدرسة، العنوان، واسم المديرة) و\textbf{تطبع} رسالة ملائمة.
\end{enumerate}

\ifwithsols
\begin{boxSolution}[البندان أ و ب]
\textbf{أ.} الإجابة الصحيحة: \textenglish{a}. \quad \textbf{ب.} الإجابة الصحيحة: \textenglish{a}.
\end{boxSolution}

\begin{boxSolution}[البند ج]
\begin{english}
\begin{minted}{python}
def school_info():
    school_name = input("Enter school name: ")
    address = input("Enter address: ")
    principal = input("Enter principal name: ")
    print("School:", school_name)
    print("Address:", address)
    print("Principal:", principal)

# Call the function
school_info()
\end{minted}
\end{english}
\end{boxSolution}

\fi

\end{enumerate}

\clearpage
%==============================================================
\section*{القسم الثاني: كتابة عمليات تتلقى معطيات}

\begin{enumerate}[resume, label=\textbf{تمرين \arabic*:}, wide=0pt, leftmargin=*, itemsep=2.2em]

% ---------------------------------------------------------------- 4
\item اكتبوا عملية باسم \textenglish{rectangle()} \textbf{تتلقى} معطيين صحيحين \textenglish{a} و \textenglish{b} يمثلان طولي ضلعين متجاورين في مستطيل. على العملية أن \textbf{تحسب} و\textbf{تطبع} \textbf{مساحة} المستطيل.

\ifwithsols
\begin{boxSolution}
\begin{english}
\begin{minted}{python}
def rectangle(a, b):
    area = a * b
    print("The area of the rectangle is:", area)

# Call the function
rectangle(4, 6)
\end{minted}
\end{english}
\end{boxSolution}

\filbreak
\fi

% ---------------------------------------------------------------- 5
\item اكتبوا عملية باسم \textenglish{pool()} \textbf{تتلقى} معطيين: عدد المرات التي نريد الذهاب بها إلى البركة (\textenglish{times}) وسعر تذكرة الدخول (\textenglish{price}). تحسب العملية و\textbf{تطبع} المبلغ \textbf{الكلي} الذي يجب دفعه.

% TODO: verify this sentence and its numbers against the original PDF (the extracted text was garbled).
\textbf{احسبوا بدون برمجة:} إذا كان سعر الاشتراك للبركة 200 ش.ج، وسعر تذكرة الدخول للمرة الواحدة 20 ش.ج، فهل من الأفضل شراء اشتراك أم دفع ثمن تذكرة في كل مرة؟

\ifwithsols
\begin{boxSolution}
\begin{english}
\begin{minted}{python}
def pool(times, price):
    total = times * price
    print("The total price is:", total)

# Call the function
pool(3, 20)
\end{minted}
\end{english}
\end{boxSolution}

\filbreak
\fi

% ---------------------------------------------------------------- 6
\item اكتبوا عملية باسم \textenglish{super\_market()} \textbf{تتلقى} ثلاثة معطيات:

\begin{itemize}[itemsep=0.3em]
\item \textenglish{juice\_bottle}: عدد قناني العصير.
\item \textenglish{bottle\_price}: سعر القنينة الواحدة.
\item \textenglish{discount}: التخفيض. (افترضوا أنّ التخفيض نسبة مئوية، مثلًا \textenglish{20} تعني \textenglish{20\%}.)
\end{itemize}

تحسب العملية و\textbf{تطبع} السعر الذي يجب دفعه لشراء قناني العصير \textbf{قبل} التخفيض و\textbf{بعده}، مع رسالة ملائمة لكل سعر.

\ifwithsols
\begin{boxSolution}
\begin{english}
\begin{minted}{python}
def super_market(juice_bottle, bottle_price, discount):
    price_before = juice_bottle * bottle_price
    price_after = price_before - price_before * discount / 100
    print("Price before discount:", price_before)
    print("Price after discount:", price_after)

# Call the function
super_market(10, 5, 20)
\end{minted}
\end{english}
\end{boxSolution}

\filbreak
\fi

% ---------------------------------------------------------------- 7
\item اكتبوا عملية باسم \textenglish{ice\_cream()} \textbf{تتلقى} معطيين: طعم البوظة وتقييمها (عدد صحيح من 1 إلى 5). \textbf{تطبع} العملية رسالة تعرض طعم البوظة وتقييمها.

استدعوا العملية \textbf{3 مرات} مع قيم مختلفة للمعطيات.

\ifwithsols
\begin{boxSolution}
\begin{english}
\begin{minted}{python}
def ice_cream(flavor, rating):
    print("The", flavor, "ice cream got a rating of", rating, "out of 5")

# Call the function
ice_cream("Vanilla", 5)
ice_cream("Chocolate", 4)
ice_cream("Pistachio", 3)
\end{minted}
\end{english}
\end{boxSolution}

\fi

\filbreak
% ---------------------------------------------------------------- 8
\item اكتبوا عملية باسم \textenglish{certificate()} \textbf{تتلقى} المعطيين التاليين:

\begin{itemize}[itemsep=0.3em]
\item \textenglish{student\_name}: اسم الطالب/ة.
\item \textenglish{message}: رسالة.
\end{itemize}

على العملية أن \textbf{تطبع} شهادة إنهاء للطالب تهنّئه على نجاحه في دراسة موضوع ``العمليات''.

\ifdetailed
\begin{boxExample}
\begin{english}
\begin{minted}{text}
*****************************
**** Well Done! ****
**** Alice ****
**** you know functions! ****
*****************************
\end{minted}
\end{english}
\end{boxExample}
\fi

\ifwithsols
\begin{boxSolution}
\begin{english}
\begin{minted}{python}
def certificate(student_name, message):
    print("*****************************")
    print("**** Well Done! ****")
    print("****", student_name, "****")
    print("****", message, "****")
    print("*****************************")

# Call the function
certificate("Alice", "you know functions!")
\end{minted}
\end{english}
\end{boxSolution}
\fi

\end{enumerate}

\vspace{1cm}
\begin{flushleft}
أرجو لكم وقتًا ممتعًا.

الأستاذ محمود اغبارية.
\end{flushleft}

\end{document}
